\chapter{Présentation du dataset}

\section{Source et volumétrie}

Le dataset utilisé est \emph{AI Jobs Market 2025-2026 | Salaries}
(Kaggle), déposé tel quel dans \texttt{data/raw/} (gitignored, jamais
commité — voir \texttt{CLAUDE.md}). Vérifié directement sur le fichier
réel :

\begin{center}
\begin{tabular}{ll}
\toprule
Caractéristique & Valeur \\
\midrule
Nombre de lignes de données & 1500 \\
Nombre de colonnes & 25 \\
Séparateur & virgule (\texttt{,}) \\
En-tête & présent (1 ligne, noms de colonnes explicites) \\
Guillemets/virgules internes & aucun — un \texttt{SerDe} délimité simple suffit \\
\bottomrule
\end{tabular}
\end{center}

\begin{constat}[Écart avec l'énoncé initial du sujet]
Le ticket \#111 (comme certaines versions de suivi antérieures)
mentionnait 26 colonnes par anticipation, avant inspection réelle du
fichier. Le comptage effectif sur le CSV (\texttt{head -1 | tr ',' '\textbackslash n' | wc -l})
donne \textbf{25 colonnes}, cohérent avec \texttt{CLAUDE.md} et le
\texttt{StructType} de \texttt{src/etl/schema.py}. Conformément à la
remarque de l'énoncé (\og les noms exacts des colonnes devront être
adaptés après inspection du fichier Kaggle, vous devez documenter toute
adaptation \fg{}), c'est le chiffre réel qui fait foi ici.
\end{constat}

\section{Variables principales}

\begin{center}
\small
\begin{tabular}{ll}
\toprule
Catégorie & Colonnes \\
\midrule
Identifiant / métier & \texttt{job\_id}, \texttt{job\_title}, \texttt{job\_category} \\
Expérience & \texttt{experience\_level}, \texttt{years\_of\_experience}, \texttt{is\_senior} \\
Formation & \texttt{education\_required} \\
Rémunération & \texttt{annual\_salary\_usd}, \texttt{salary\_min\_usd}, \texttt{salary\_max\_usd}, \texttt{salary\_tier} \\
Localisation & \texttt{city}, \texttt{country} \\
Organisation & \texttt{company\_size}, \texttt{industry} \\
Travail à distance & \texttt{remote\_work}, \texttt{is\_remote\_friendly} \\
Compétences & \texttt{required\_skills} (liste \texttt{\textbar}-séparée) \\
Marché IA & \texttt{ai\_salary\_premium\_pct}, \texttt{demand\_score}, \texttt{demand\_growth\_yoy\_pct}, \texttt{is\_llm\_role} \\
Temporel & \texttt{posting\_year}, \texttt{posting\_month} \\
Autres & \texttt{benefits\_score\_10} \\
\bottomrule
\end{tabular}
\end{center}

Ces colonnes couvrent directement les axes demandés par l'énoncé :
offres, métiers, compétences, salaires, pays, secteurs et dates.
\texttt{required\_skills} est la seule colonne multi-valuée (transformée
en table \texttt{competences} au chapitre~7 via \texttt{split}/\texttt{explode}) ;
toutes les autres colonnes sont directement exploitables telles quelles.

\section{Premières observations de qualité (avant nettoyage)}

Contrôle exécuté sur le fichier brut, avant toute transformation Spark
(distinct des comptages post-nettoyage du chapitre~6, qui portent sur
le DataFrame après lecture HDFS) :

\begin{center}
\begin{tabular}{ll}
\toprule
Contrôle & Résultat \\
\midrule
Valeurs manquantes (25 colonnes, 1500 lignes) & 0 \\
\texttt{job\_id} en double & 0 \\
\texttt{annual\_salary\_usd} $\leq$ 0 & 0 \\
Métiers distincts (\texttt{job\_title}) & 25 \\
Pays distincts (\texttt{country}) & 14 \\
Salaire annuel (\texttt{annual\_salary\_usd}) & min 90\,000\$ / max 384\,000\$ / moyenne 194\,892\$ \\
Répartition par expérience & Entry 385, Mid 370, Senior 364, Lead 381 \\
\bottomrule
\end{tabular}
\end{center}

\begin{constat}[Un dataset propre par construction]
Aucune anomalie détectée sur les critères contrôlés (NULL, doublons,
salaires invalides) — cohérent avec le résultat obtenu après nettoyage
Spark au chapitre~6 (0 ligne affectée par les règles de nettoyage). Ce
dataset est vraisemblablement généré/synthétisé plutôt que collecté
directement (répartition d'expérience quasi uniforme, absence totale de
valeur manquante sur 25 colonnes) — un point repris au chapitre~13
(limites de l'étude) : la propreté du dataset ne garantit pas sa
représentativité du marché réel de l'emploi IA.
\end{constat}
